\chapter{Functions}
\label{ch:functions}

Functions are the primary unit of code organization in \haira{}. They are first-class values that can be passed as arguments, returned from other functions, and stored in variables.

\section{Function Declarations}

\begin{lstlisting}
fn function_name(param1: Type1, param2: Type2) -> ReturnType {
    // body
}
\end{lstlisting}

\subsection{Basic Examples}

\begin{lstlisting}
fn greet() {
    io.println("Hello!")
}

fn add(a: int, b: int) -> int {
    return a + b
}

fn is_positive(n: int) -> bool {
    return n > 0
}
\end{lstlisting}

\subsection{No Return Type}

Functions without a return type implicitly return \code{()}:

\begin{lstlisting}
fn log_message(msg: string) {
    io.println("[LOG] " + msg)
}
\end{lstlisting}

\section{Parameters}

\subsection{Required Parameters}

All parameters must have explicit type annotations:

\begin{lstlisting}
fn calculate(x: int, y: int, z: float) -> float {
    return (x + y) as float * z
}
\end{lstlisting}

\subsection{Default Parameters}

Parameters can have default values:

\begin{lstlisting}
fn greet(name: string, greeting: string = "Hello") {
    io.println(greeting + ", " + name + "!")
}

greet("Alice")              // "Hello, Alice!"
greet("Bob", "Hi")          // "Hi, Bob!"
\end{lstlisting}

\begin{note}
Parameters with defaults must come after parameters without defaults.
\end{note}

\subsection{Named Arguments}

Functions can be called with named arguments:

\begin{lstlisting}
fn create_user(name: string, email: string, age: int = 0) -> User {
    return User{name: name, email: email, age: age}
}

// Positional
create_user("Alice", "alice@example.com", 30)

// Named
create_user(name: "Bob", email: "bob@example.com")

// Mixed (positional first, then named)
create_user("Charlie", email: "charlie@example.com", age: 25)
\end{lstlisting}

\subsection{Variadic Parameters}

Functions can accept variable numbers of arguments:

\begin{lstlisting}
fn sum(numbers: ...int) -> int {
    total = 0
    for n in numbers {
        total += n
    }
    return total
}

sum(1, 2, 3)        // 6
sum(1, 2, 3, 4, 5)  // 15
\end{lstlisting}

\section{Return Values}

\subsection{Single Return}

\begin{lstlisting}
fn double(n: int) -> int {
    return n * 2
}
\end{lstlisting}

\subsection{Multiple Returns}

Functions can return multiple values using tuples:

\begin{lstlisting}
fn divide(a: int, b: int) -> (int, int) {
    return a / b, a % b
}

quotient, remainder = divide(17, 5)  // 3, 2
\end{lstlisting}

\subsection{Implicit Return}

The last expression in a function body is implicitly returned:

\begin{lstlisting}
fn add(a: int, b: int) -> int {
    a + b  // Implicit return
}

fn max(a: int, b: int) -> int {
    if a > b { a } else { b }  // Implicit return
}
\end{lstlisting}

\subsection{Early Return}

Use \keyword{return} for early exit:

\begin{lstlisting}
fn factorial(n: int) -> int {
    if n <= 1 {
        return 1
    }
    return n * factorial(n - 1)
}
\end{lstlisting}

\section{Function Types}

Functions have types that can be used in type annotations:

\begin{lstlisting}
// Function type: fn(int, int) -> int
type BinaryOp = fn(int, int) -> int

fn apply(op: BinaryOp, a: int, b: int) -> int {
    return op(a, b)
}

fn add(a: int, b: int) -> int { a + b }
fn mul(a: int, b: int) -> int { a * b }

apply(add, 2, 3)  // 5
apply(mul, 2, 3)  // 6
\end{lstlisting}

\section{Closures}

Closures are anonymous functions that can capture values from their environment:

\begin{lstlisting}
fn make_counter() -> fn() -> int {
    count = 0
    return fn() -> int {
        count += 1
        return count
    }
}

counter = make_counter()
counter()  // 1
counter()  // 2
counter()  // 3
\end{lstlisting}

\subsection{Closure Syntax}

\begin{lstlisting}
// Full syntax
fn(a: int, b: int) -> int {
    return a + b
}

// Short syntax (single expression)
fn(a: int, b: int) -> int { a + b }

// With type inference (in context)
numbers = [1, 2, 3]
doubled = array.map(numbers, fn(n) { n * 2 })
\end{lstlisting}

\subsection{Capture Semantics}

Closures capture variables by reference:

\begin{lstlisting}
fn example() {
    values = []

    for i in 0..3 {
        // Captures 'i' by reference
        values = array.push(values, fn() { i })
    }

    // All closures see final value of i
    values[0]()  // 3
    values[1]()  // 3
    values[2]()  // 3
}
\end{lstlisting}

To capture by value, use explicit binding:

\begin{lstlisting}
fn example() {
    values = []

    for i in 0..3 {
        captured = i  // Capture current value
        values = array.push(values, fn() { captured })
    }

    values[0]()  // 0
    values[1]()  // 1
    values[2]()  // 2
}
\end{lstlisting}

\section{Methods}

Methods are functions associated with a type:

\begin{lstlisting}
struct Rectangle {
    width: float
    height: float
}

Rectangle.area() -> float {
    return self.width * self.height
}

Rectangle.perimeter() -> float {
    return 2.0 * (self.width + self.height)
}

Rectangle.scale(factor: float) -> Rectangle {
    return Rectangle{
        width = self.width * factor,
        height = self.height * factor
    }
}

// Usage
rect = Rectangle{width = 10.0, height = 5.0}
rect.area()       // 50.0
rect.perimeter()  // 30.0
rect.scale(2.0)   // Rectangle{width: 20.0, height: 10.0}
\end{lstlisting}

\subsection{Method Semantics}

Methods use \code{self} to access the receiver. All methods can mutate the receiver (pointer semantics). Method visibility is inherited from the type: if the type is \keyword{pub}, all its methods are accessible to importers. Methods cannot be defined on imported types.

\begin{lstlisting}
struct Point {
    x: float
    y: float
}

Point.distance_to(other: Point) -> float {
    dx = self.x - other.x
    dy = self.y - other.y
    return math.sqrt(dx*dx + dy*dy)
}

Point.translate(dx: float, dy: float) {
    self.x += dx
    self.y += dy
}
\end{lstlisting}

\section{Generic Functions}

Functions can be parameterized over types:

\begin{lstlisting}
fn identity<T>(value: T) -> T {
    return value
}

fn swap<T, U>(pair: (T, U)) -> (U, T) {
    return (pair.1, pair.0)
}

fn first<T>(items: []T) -> T? {
    if array.len(items) == 0 {
        return nil
    }
    return items[0]
}
\end{lstlisting}

\subsection{Type Constraints}

Generic types can be constrained:

\begin{lstlisting}
constraint Comparable {
    fn compare(other: Self) -> int
}

fn max<T: Comparable>(a: T, b: T) -> T {
    if a.compare(b) > 0 {
        return a
    }
    return b
}

fn sort<T: Comparable>(items: []T) -> []T {
    // Sorting implementation
}
\end{lstlisting}

\section{Higher-Order Functions}

Functions that take or return other functions:

\begin{lstlisting}
// Function as parameter
fn apply_twice(f: fn(int) -> int, x: int) -> int {
    return f(f(x))
}

fn double(n: int) -> int { n * 2 }

apply_twice(double, 5)  // 20

// Function as return value
fn make_multiplier(factor: int) -> fn(int) -> int {
    return fn(n: int) -> int {
        return n * factor
    }
}

triple = make_multiplier(3)
triple(7)  // 21
\end{lstlisting}

\section{Recursion}

Functions can call themselves:

\begin{lstlisting}
fn factorial(n: int) -> int {
    if n <= 1 {
        return 1
    }
    return n * factorial(n - 1)
}

fn fibonacci(n: int) -> int {
    match n {
        0 => 0
        1 => 1
        _ => fibonacci(n - 1) + fibonacci(n - 2)
    }
}
\end{lstlisting}

\subsection{Tail Call Optimization}

\haira{} optimizes tail-recursive functions:

\begin{lstlisting}
fn factorial_tail(n: int, acc: int = 1) -> int {
    if n <= 1 {
        return acc
    }
    return factorial_tail(n - 1, n * acc)  // Tail call
}
\end{lstlisting}

\section{Function Overloading}

\haira{} does not support function overloading. Each function name must be unique within its scope. Use different names or generic functions instead:

\begin{lstlisting}
// Instead of overloading, use descriptive names
fn add_ints(a: int, b: int) -> int { a + b }
fn add_floats(a: float, b: float) -> float { a + b }

// Or use generics
fn add<T: Addable>(a: T, b: T) -> T { a + b }
\end{lstlisting}

\section{The main Function}

Every executable \haira{} program must have a \code{main} function:

\begin{lstlisting}
fn main() {
    io.println("Hello, World!")
}

// With exit code
fn main() -> int {
    io.println("Hello, World!")
    return 0
}
\end{lstlisting}

\section{Function Grammar}

\begin{lstlisting}[language=EBNF]
function_decl = "fn" [ receiver ] identifier [ type_params ]
                "(" [ params ] ")" [ "->" type ] block ;

receiver = "(" identifier ":" [ "*" ] type ")" ;

type_params = "<" type_param { "," type_param } ">" ;

type_param = identifier [ ":" constraint_list ] ;

constraint_list = identifier { "+" identifier } ;

params = param { "," param } ;

param = identifier ":" [ "..." ] type [ "=" expression ] ;

closure = "fn" "(" [ closure_params ] ")" [ "->" type ]
          ( block | expression ) ;

closure_params = closure_param { "," closure_param } ;

closure_param = identifier [ ":" type ] ;
\end{lstlisting}
